% AI-written, AI cross-reviewed; no human sign-off; not a claim that RH is proved
% RH Notes blueprint source excerpts for Z1.5.07
% Verbatim excerpts in the notation of the parent blueprint. The labelled editions and source ranges appear below.
% See the unit YAML for the current proof status. Earlier-round passages retain their original wording.
\section*{Z1.5.07. CM Frobenius trace bound}

% Parent source SHA-256: 1af64922b98ed2310744411b1f18b25cc426c97220646321471836e1a3a409d6
% blueprint/ch_P.tex lines 12504 through 12520; r5:small:Z1.5.07
\subsection{Z1.5.07. Turning a rank-three CM classifier into a trace classifier}
\label{r5:small:Z1.5.07}
\textit{Complete under $\mathrm{CountKAlt}$.} Fix a CM field $K$. The constants work for every $e,N\ge1$, $p$ prime away from $N$ and the discriminant of $K$, and every $w\mid p$. Write $q=\mathrm{N}w$. For one finite trace datum, chosen before the odd coefficient prime $\ell\ne p$, every $\mathrm{Adm3NAlt}$ representation at every place above $w$ has a label. Equal labels imply isomorphic representations and that label evaluates to the trace at \emph{inverse} Frobenius. Its cardinality is bounded by $c'[N(e+1)]^A$. The two numerical conditions on this trace datum are $|t_s|\le\sqrt q+2$ and integrality of $q^e t_s$.
\par{\footnotesize \nolinkurl{UpperTheory.GalN.V.Win.N.frobenius_trace_count_bound_of_cm_eigensystem_count_alt} (\nolinkurl{Development/AutomorphicArithmetic/AutomorphicVarieties/ModularityArgument/ThresholdArguments/DenseThresholdCases/FrobeniusTraceCountBoundOfCmEigensystemCount.lean}\allowbreak:418--461). \nolinkurl{UpperTheory.GalN.V.Win.TrData} (\nolinkurl{Development/AutomorphicArithmetic/Foundations/GaloisRepresentations/AdmissibleResidues/WinDefs.lean}\allowbreak:70--78). \nolinkurl{UpperTheory.GalN.V.Win.TrData.Classifies} (\nolinkurl{Development/AutomorphicArithmetic/Foundations/GaloisRepresentations/AdmissibleResidues/WinDefs.lean}\allowbreak:87--90).\par}

Apply the assumed rank-three count with $m=(e-1,-1,-e-1)$, a strictly decreasing tuple paired to sum $-2$, and $H=e+1$. If $\beta_s$ is its multiset of three roots, keep its index set and put
\[
           t_s=\sum_{b\in\beta_s}(qb)^{-1}.
\]
The label permutations under field automorphisms still work since the operations commute with them. Let $a=\sqrt q\ge1$ and let $x,y,z$ be the three numbers $|(qb)^{-1}|$. The window inequalities and the product clause give $a^{-1}\le x,y,z\le a$ and $xyz=1$. Two, say $x,y$, are on the same side of one. Thus $x+y\le 1+xy=1+1/z$. For $1/a\le z\le a$, $(a-z)(az-1)\ge0$ implies $z+1/z\le a+1/a\le a+1$. This proves the required bound by the triangle inequality. Since all three weights are $\le e-1$, the inverse-integrality clause makes $q^{e-1}b^{-1}$ integral, and summing gives integrality of $q^e t_s$. The identification of label sets preserves cardinality. The original $c[N(e+2)]^A$ is at most $c2^A[N(e+1)]^A$.
\par{\footnotesize \nolinkurl{UpperTheory.GalN.V.Win.traces_from_eigenvalues_rank3} (\nolinkurl{Development/AutomorphicArithmetic/Foundations/GaloisRepresentations/AdmissibleResidues/FrobeniusTraceCountBoundOfRatEigensystemCount.lean}\allowbreak:78--145). \nolinkurl{UpperTheory.GalN.V.Win.N.frobenius_trace_count_bound_of_cm_eigensystem_count_alt} (\nolinkurl{Development/AutomorphicArithmetic/AutomorphicVarieties/ModularityArgument/ThresholdArguments/DenseThresholdCases/FrobeniusTraceCountBoundOfCmEigensystemCount.lean}\allowbreak:421--448).\par}

For the actual representation $r$ label $r\otimes\chi_{\rm cyc}^{-1}$. Its automorphic condition is already the one in the starting record. Twisting preserves invariant subspaces and shifts all ordinary weights by $-1$. Away from $\ell$, inertia acts trivially on the cyclotomic twist. Above $\ell$, twist every diagonal character in the same triangular basis. Equality of two characters after the common twist cancels; in the potentially crystalline branch use the fixed local bridge and the period comparison in Section~\ref{r6:review:Z1.5.07}. The oddness clause is similarly preserved under a common character twist. Equal labels give a conjugating matrix for the twisted representations and the scalar twist cancels in the conjugacy identity. Since $\ell\ne p$, B\'ezout for $p,\ell$ excludes $\ell$ from the ideal $w$. At the given Frobenius $\sigma$ the cyclotomic character is $q$, so at $\sigma^{-1}$ the inverse cyclotomic twist multiplies $r(\sigma^{-1})$ by $q$. The classifier gives the multiset of inverse roots with multiplicity for the twisted representation at inverse Frobenius, whence $q\,\iota\operatorname{tr}r(\sigma^{-1})=\sum_{b\in\beta_s}b^{-1}$. Dividing by $q$ gives precisely $t_s$.
\par{\footnotesize \nolinkurl{UpperTheory.GalN.V.Win.N.admissible_twist_cyclotomic} (\nolinkurl{Development/AutomorphicArithmetic/AutomorphicVarieties/ModularityArgument/ThresholdArguments/DenseThresholdCases/FrobeniusTraceCountBoundOfCmEigensystemCount.lean}\allowbreak:409--417). \nolinkurl{UpperTheory.level_twist_cyclotomic_pow} (\nolinkurl{Development/AutomorphicArithmetic/AutomorphicVarieties/BuildingsResidueChars/LevelTwistCyclotomicPow.lean}\allowbreak:19--59). \nolinkurl{UpperTheory.map_trace_eq_sum_roots_of_charpoly_map} (\nolinkurl{Development/AutomorphicArithmetic/AutomorphicVarieties/BuildingsResidueChars/MapTraceEqSumRootsOfCharpolyMap.lean}\allowbreak:5--49). \nolinkurl{ReciprocityTheory.Gal.cycZ_of_isFrob} (\nolinkurl{ReciprocityTheory/GaloisValuations/Global.lean}\allowbreak:364--380). \nolinkurl{UpperTheory.GalN.V.Win.N.frobenius_trace_count_bound_of_cm_eigensystem_count_alt} (\nolinkurl{Development/AutomorphicArithmetic/AutomorphicVarieties/ModularityArgument/ThresholdArguments/DenseThresholdCases/FrobeniusTraceCountBoundOfCmEigensystemCount.lean}\allowbreak:447--461).\par}

\textit{[Hypothesis, discharged at the use:]} $\mathrm{CountKAlt}$ is Z1.5.03. The consumer passes the theorem obtained from its given automorphic and level inputs to this trace conversion. The proof for a fixed CM field above is conditional on precisely that count.
\par{\footnotesize \nolinkurl{UpperTheory.GalN.V.Win.N.countK_alt} (\nolinkurl{Development/AutomorphicArithmetic/AutomorphicVarieties/ModularityArgument/ThresholdArguments/DenseThresholdCases/FrobeniusTraceCountBoundOfCmEigensystemCount.lean}\allowbreak:63--68). \nolinkurl{UpperTheory.GalN.V.Win.N.numberU3U_of_adm2A_imp_adm2B} (\nolinkurl{Development/AutomorphicArithmetic/AutomorphicVarieties/ModularityArgument/ThresholdArguments/DenseThresholdCases/KeepNumU.lean}\allowbreak:25--32).\par}

% blueprint/ch_P.tex lines 12587 through 12593; r6:review:Z1.5.07
\subsection{Z1.5.07. The potentially crystalline alternative of the twist}\label{r6:review:Z1.5.07} \emph{Original argument on p.~\pageref{r5:small:Z1.5.07}.} 

Above $\ell$ the level hypothesis permits either a separation of diagonal characters or a potentially crystalline model. In the second case keep the bridge $B_w$ and a finite local extension $K'$ where the supplied model of $r$ is crystalline. Enlarge its coefficient field and that of the one-dimensional character to their compositum. Extension of coefficients keeps the crystalline model over $K'$; scalar multiplication by the rank-one model gives a model of the twist, since the two original model identities multiply entry by entry. On $G_{K'}$ its character is the $j$-th power of the local cyclotomic character. This identification follows from the bridge identity for the cyclotomic action, and applies to $j=-1$ here.
\par{\footnotesize \nolinkurl{UpperTheory.level_twist_cyclotomic_pow} (\nolinkurl{Development/AutomorphicArithmetic/AutomorphicVarieties/BuildingsResidueChars/LevelTwistCyclotomicPow.lean}\allowbreak:19--61). \nolinkurl{GlobalTheory.HodgeOfTrue.potCrys_twistBy} (\nolinkurl{Development/AbelianVarietyTheory/Proofs/InertiaEigenvalues/GaloisHodgeOfCuspidal.lean}\allowbreak:188--234). \nolinkurl{UpperTheory.isCrysAt_cyclotomic_zpow_restrict_decompositionGroup} (\nolinkurl{Development/AutomorphicArithmetic/AutomorphicVarieties/BuildingsResidueChars/IsCrysAtCyclotomicZpowRestrictDecompositionGroup.lean}\allowbreak:7--43).\par}

Write $B$ for the crystalline period algebra of the given data, with its unit $t$ satisfying $g(t)=\chi_{\rm cyc}(g)t$. On the tensor product of the coefficient module with $B$, multiplication on the period factor by $t^j$ takes the twisted action to the original action: acting on $t^j b$ contributes exactly $\chi_{\rm cyc}(g)^j$ to $g(b)$. Its inverse is multiplication by $t^{-j}$. It restricts to an isomorphism of the two modules of invariants over the fixed coefficient scalars. Hence the invariant dimension for the twisted model over $K'$ still has its full crystalline value. This proves the required potential crystallinity. For the rank-one character the same computation with the trivial representation gives its crystallinity; the general twist lemma in the source also permits an arbitrary crystalline character using tensor and dual closure. Continuity of the twist follows by multiplying the continuous character and representation. Finally the classifier already evaluates the characteristic polynomial of the twist at $\sigma^{-1}$ with roots $b^{-1}$; division of its trace by $q$ gives the required value. This deduction takes $\mathrm{CountKAlt}$ as its counting input.
\par{\footnotesize \nolinkurl{GlobalTheory.BlockCrys.isCrys_mul_cyc}, \nolinkurl{isCrys_cyc_zpow} (\nolinkurl{Development/AbelianVarietyTheory/Proofs/InertiaEigenvalues/BlockCharactersCrystallineOfRules.lean}\allowbreak:28--72). \nolinkurl{GlobalTheory.TopSteps.CrysRules.isCrys_iff} (\nolinkurl{Development/AbelianVarietyTheory/Proofs/InertiaEigenvalues/CrystallineSumsTwistsCharactersExpectedJumps.lean}\allowbreak:129--136). \nolinkurl{UpperTheory.GalN.V.Win.EigData.Classifies} (\nolinkurl{Development/AutomorphicArithmetic/Foundations/GaloisRepresentations/AdmissibleResidues/WinDefs.lean}\allowbreak:34--38). \nolinkurl{UpperTheory.GalN.V.Win.N.admissible_twist_cyclotomic} (\nolinkurl{Development/AutomorphicArithmetic/AutomorphicVarieties/ModularityArgument/ThresholdArguments/DenseThresholdCases/FrobeniusTraceCountBoundOfCmEigensystemCount.lean}\allowbreak:409--417). \nolinkurl{UpperTheory.GalN.V.Win.N.frobenius_trace_count_bound_of_cm_eigensystem_count_alt} (\nolinkurl{Development/AutomorphicArithmetic/AutomorphicVarieties/ModularityArgument/ThresholdArguments/DenseThresholdCases/FrobeniusTraceCountBoundOfCmEigensystemCount.lean}\allowbreak:418--461).\par}

% Parent source SHA-256: 1fdc6ea9aac6cf10316f85b3253b33ae98289a76ff5aa2e79daaaf11e057ce18
% blueprint/ch_P_r7.tex lines 23 through 24; r7:review:Z1.5.07
\subsection{The crystalline twist in Z1.5.07}\label{r7:review:Z1.5.07}
For the conditional trace classifier the assumption is $\mathrm{CountKAlt}$. In its potentially crystalline branch choose the same local period comparison and a common finite field of definition as on page~\pageref{r6:review:Z1.5.07}. Multiplication by $t^j$ intertwines the cyclotomically twisted action with the original one since $g(t)=\chi_{\mathrm{cyc}}(g)t$; $t^{-j}$ gives the inverse. Thus the invariant dimensions agree. The general source theorem allows a crystalline character and takes a common finite extension, while this classifier needs only $j=-1$. The local representation remains continuous since scalar multiplication and the two representations are continuous. At inverse Frobenius the classifier for the twist already gives the inverse roots. Undoing the twist divides these roots by the residue norm, as in \nolinkurl{admissible_twist_cyclotomic}, \nolinkurl{frobenius_trace_count_bound_of_cm_eigensystem_count_alt} (\path{Development/AutomorphicArithmetic/AutomorphicVarieties/ModularityArgument/ThresholdArguments/DenseThresholdCases/FrobeniusTraceCountBoundOfCmEigensystemCount.lean:409--461}). The period calculation is \nolinkurl{GlobalTheory.BlockCrys.isCrys_mul_cyc}, \nolinkurl{isCrys_cyc_zpow} (\path{Development/AbelianVarietyTheory/Proofs/InertiaEigenvalues/BlockCharactersCrystallineOfRules.lean:28--72}); the extension and tensor argument is \nolinkurl{GlobalTheory.HodgeOfTrue.potCrys_twistBy} (\path{Development/AbelianVarietyTheory/Proofs/InertiaEigenvalues/GaloisHodgeOfCuspidal.lean:188--234}). The unit supplies the implication with the original count as premise.

